\documentclass[aspectratio=169]{beamer}
\input{header}
\coursecode{JKSSB}

\title{Keyboard Shortcuts}
\subtitle{Lesson 21 --- High-Yield Keys and Combinations}
\author{JKSSB Computer Awareness}
\date{\today}

\begin{document}
\frame{\titlepage}

\begin{frame}{Reading a Keyboard Shortcut}
  \begin{itemize}
    \item A \key{keyboard shortcut} is a key or a combination of keys that
      performs a command directly, without opening a menu.
    \item In a combination such as \key{Ctrl+C}, the plus sign means: hold the
      first key down, then tap the second.
    \item The modifier keys are \key{Ctrl} (Control), \key{Alt} (Alternate),
      \key{Shift}, and the \key{Windows (Win)} key.
    \item \key{Function keys} \texttt{F1}--\texttt{F12} form the top row of
      the keyboard.
  \end{itemize}
  \begin{block}{The one idea}
    A shortcut is the same command as the menu item --- just reached with the
    keyboard instead of the mouse.
  \end{block}
\end{frame}

\begin{frame}{The Clipboard Trio: Copy, Cut, Paste}
  \begin{center}
    \begin{tabular}{ll}
      \toprule
      \textbf{Shortcut} & \textbf{Action} \\
      \midrule
      \key{Ctrl+C} & Copy the selected item \\
      \key{Ctrl+X} & Cut the selected item \\
      \key{Ctrl+V} & Paste the copied or cut item \\
      \bottomrule
    \end{tabular}
  \end{center}
  \vspace{0.25cm}
  \begin{alertblock}{Trap alert}
    \key{Copy} leaves the original in place; \key{Cut} removes the original.
    Do not treat \texttt{Ctrl+C} and \texttt{Ctrl+X} as the same command.
  \end{alertblock}
\end{frame}

\begin{frame}{Undo, Redo, Select All}
  \begin{itemize}
    \item \key{Ctrl+Z} --- Undo the last action.
    \item \key{Ctrl+Y} --- Redo the action that was undone.
    \item \key{Ctrl+A} --- Select All the content in the current window or
      document.
  \end{itemize}
  \begin{block}{Undo and redo are a pair}
    \texttt{Ctrl+Z} steps backward; \texttt{Ctrl+Y} steps forward again. They
    are not interchangeable.
  \end{block}
\end{frame}

\begin{frame}{Save, Print, Find, Replace}
  \begin{center}
    \begin{tabular}{ll}
      \toprule
      \textbf{Shortcut} & \textbf{Action} \\
      \midrule
      \key{Ctrl+S} & Save the current file \\
      \key{Ctrl+P} & Print \\[0.2em]
      \key{Ctrl+F} & Find (search) text \\
      \key{Ctrl+H} & Replace text \\
      \bottomrule
    \end{tabular}
  \end{center}
  \vspace{0.2cm}
  \muted{Find locates a word; Replace finds it and substitutes another.}
\end{frame}

\begin{frame}{Trap Alert: The Ctrl Near-Neighbours}
  \begin{itemize}
    \item \key{Ctrl+Z} = Undo \hfill \muted{not} \hfill \key{Ctrl+Y} = Redo.
    \item \key{Ctrl+F} = Find \hfill \muted{not} \hfill \key{Ctrl+H} =
      Replace.
    \item \key{Ctrl+C} = Copy \hfill \muted{not} \hfill \key{Ctrl+X} = Cut.
    \item \key{Ctrl+S} = Save \hfill \muted{not} \hfill \key{Ctrl+P} = Print.
  \end{itemize}
  \begin{alertblock}{Read the letter, not the habit}
    The whole mark of these items is the \key{second key}. Slow down and read
    the exact letter.
  \end{alertblock}
\end{frame}

\begin{frame}{Alt-Key Combinations}
  \begin{itemize}
    \item \key{Alt+Tab} --- Switch between the open windows or running
      applications.
    \item \key{Alt+F4} --- Close the active window or application.
    \item On the desktop with no window open, \texttt{Alt+F4} opens the
      shutdown / sign-out dialog.
  \end{itemize}
  \begin{block}{Alt means an alternative action}
    \texttt{Alt} combinations mostly manage \emph{windows}, while
    \texttt{Ctrl} combinations manage \emph{content}.
  \end{block}
\end{frame}

\begin{frame}{Trap Alert: Alt+Tab vs Alt+F4}
  \begin{center}
    \begin{tabular}{ll}
      \toprule
      \textbf{Shortcut} & \textbf{What it does} \\
      \midrule
      \key{Alt+Tab} & Moves to another open window (does \emph{not} close one) \\
      \key{Alt+F4} & Closes the active window (does \emph{not} just switch) \\
      \bottomrule
    \end{tabular}
  \end{center}
  \vspace{0.25cm}
  \begin{alertblock}{Trap alert}
    \texttt{Alt+Tab} \key{switches}; \texttt{Alt+F4} \key{closes}. Reversing
    the two is a common wrong answer.
  \end{alertblock}
\end{frame}

\begin{frame}{The Function Keys}
  \begin{itemize}
    \item Function keys are \key{single keys} on the top row: \texttt{F1} to
      \texttt{F12}.
    \item Their action depends on the application that is active.
    \item The high-yield set for the exam is \key{F1}, \key{F2}, \key{F5}, and
      \key{F7}.
    \item Pressing a function key alone is not a combination; no modifier is
      held.
  \end{itemize}
  \begin{block}{One key, one job}
    Learn each high-yield function key with its standard job before learning
    application-specific behaviour.
  \end{block}
\end{frame}

\begin{frame}{Function-Key Quick Table}
  \begin{center}
    \begin{tabular}{ll}
      \toprule
      \textbf{Key} & \textbf{Standard action} \\
      \midrule
      \key{F1} & Help \\
      \key{F2} & Rename the selected file or folder \\
      \key{F5} & Refresh the window / start a slide show from the beginning \\
      \key{F7} & Spelling and grammar check (MS Word) \\
      \bottomrule
    \end{tabular}
  \end{center}
  \vspace{0.2cm}
  \muted{F1 help, F2 rename, F5 refresh/slide show, F7 spell check.}
\end{frame}

\begin{frame}{F5 and Shift+F5 in PowerPoint}
  \begin{itemize}
    \item \key{F5} --- Start the slide show \key{from the beginning}.
    \item \key{Shift+F5} --- Start the slide show \key{from the current
      slide}.
    \item In Windows (not in a slide show), \key{F5} means \key{Refresh} the
      folder or window.
  \end{itemize}
  \begin{exampleblock}{Documented exam fact}
    Website Operator 2026 asked the key that starts a slide show from the
    beginning --- the answer is \key{F5} (PYQ WO 2026 Q80).
  \end{exampleblock}
\end{frame}

\begin{frame}{The Windows-Key Shortcuts}
  \begin{center}
    \begin{tabular}{ll}
      \toprule
      \textbf{Shortcut} & \textbf{Action} \\
      \midrule
      \key{Win+D} & Show the desktop (minimise all windows) \\
      \key{Win+E} & Open File Explorer \\
      \key{Win+L} & Lock the computer \\
      \bottomrule
    \end{tabular}
  \end{center}
  \vspace{0.2cm}
  \muted{\texttt{D} for Desktop, \texttt{E} for Explorer, \texttt{L} for Lock.}
\end{frame}

\begin{frame}{The Shift Key}
  \begin{itemize}
    \item \key{Shift} produces the capital letter or the symbol on the upper
      part of a key.
    \item \key{Shift+Click} selects a \key{range} of items between two clicks.
    \item \key{Shift+F5} starts a slide show from the current slide.
    \item Held with a letter, Shift does not save, close, or print anything;
      it changes \emph{what is typed or selected}.
  \end{itemize}
\end{frame}

\begin{frame}{Key to Action: Recall Table}
  \begin{center}
    \begin{tabular}{llll}
      \toprule
      \textbf{Key} & \textbf{Action} & \textbf{Key} & \textbf{Action} \\
      \midrule
      Ctrl+C & Copy & Ctrl+S & Save \\
      Ctrl+X & Cut & Ctrl+P & Print \\
      Ctrl+V & Paste & Ctrl+F & Find \\
      Ctrl+Z & Undo & Ctrl+H & Replace \\
      Ctrl+Y & Redo & F1 & Help \\
      Ctrl+A & Select All & F2 & Rename \\
      Alt+Tab & Switch window & F5 & Refresh / slide show \\
      Alt+F4 & Close window & F7 & Spell check \\
      \bottomrule
    \end{tabular}
  \end{center}
\end{frame}

\begin{frame}{Trap Alert: Most-Missed Near-Neighbours}
  \begin{itemize}
    \item \key{F1} Help \hfill \muted{not} \hfill \key{F2} Rename.
    \item \key{F5} Refresh or slide show \hfill \muted{not} \hfill \key{F7}
      Spell check.
    \item \key{Ctrl+Z} Undo \hfill \muted{not} \hfill \key{Ctrl+Y} Redo.
    \item \key{Ctrl+F} Find \hfill \muted{not} \hfill \key{Ctrl+H} Replace.
    \item \key{Alt+Tab} Switch \hfill \muted{not} \hfill \key{Alt+F4} Close.
  \end{itemize}
\end{frame}

\begin{frame}{Lesson 21: Recall}
  \begin{itemize}
    \item Clipboard: \key{Ctrl+C} copy, \key{Ctrl+X} cut, \key{Ctrl+V} paste.
    \item Edit: \key{Ctrl+Z} undo, \key{Ctrl+Y} redo, \key{Ctrl+A} select all.
    \item File and search: \key{Ctrl+S} save, \key{Ctrl+P} print, \key{Ctrl+F}
      find, \key{Ctrl+H} replace.
    \item Windows: \key{Alt+Tab} switch, \key{Alt+F4} close, \key{Win+D} show
      desktop, \key{Win+E} File Explorer, \key{Win+L} lock.
    \item Function keys: \key{F1} help, \key{F2} rename, \key{F5}
      refresh/slide show, \key{F7} spell check.
    \item \key{Shift+F5} starts a slide show from the current slide.
  \end{itemize}
\end{frame}
\end{document}
